% Research contribution for the pinned ProveIt manuscript.
% Standard theorem environments and amsmath/amssymb are assumed.
\section{A complete formal classification of rational companion values}
\label{jnew:sec:classification}

Put
\[
 J(x)=\int_0^1\frac{\log(1+t^x)}{1+t}\,dt,
 \qquad
 F(x)=\sum_{n\ge1}\frac{\psi(nx)-\log(nx)}n.
\]
Radchenko--Zagier's identity~\cite{RadchenkoZagier} is
\begin{equation}\label{jnew:eq:JF}
 J(x)=F(2x)-2F(x)+F(x/2)+\frac{\pi^2}{12x}.
\end{equation}
All identities in this section at positive real arguments use real logarithms.
Our negative statements concern the rationalized five-term calculus, not
numerical independence of periods.

Write $A(K)=K^\times\otimes_{\mathbb Z}\mathbb Q$ additively, using
$\langle u\rangle$ for the class of $u$. Roots of unity disappear.
Let $\mathcal P(K)$ be the rationalized pre-Bloch group, with
$\partial[z]=\langle z\rangle\wedge\langle1-z\rangle$ and
$[0]=[1]=0$. For coprime positive $p,q$, let $\xi_{p/q}$ denote the
Galois-invariant element supplied by the finite cyclotomic formula for
$F(p/q)-F(1)$. Its normalization is
\begin{equation}\label{jnew:eq:xi-boundary}
 \partial\xi_{p/q}
 =\beta_q(p)-\beta_p(q)-\langle p\rangle\wedge\langle q\rangle,
 \qquad
 \beta_m(a)=\sum_{k=1}^{m-1}
 \langle1-\zeta_m^{ak}\rangle\wedge\langle1-\zeta_m^k\rangle.
\end{equation}
Set $\beta_1=\beta_2=0$. Fractions in subscripts are always reduced
before using \eqref{jnew:eq:xi-boundary}. Define
\begin{equation}\label{jnew:eq:eta}
 \eta_{p/q}=\xi_{2p/q}-2\xi_{p/q}+\xi_{p/(2q)}
 \quad\text{in }\mathcal P(\mathbb Q(\zeta_{2pq})).
\end{equation}
A formal logarithmic reduction of $J(p/q)$ means $\eta_{p/q}=0$;
a formal rational-dilogarithm reduction means
$\eta_{p/q}\in\operatorname{im}\mathcal P(\mathbb Q)$.
Products of algebraic logarithms and rational multiples of $\pi^2$
are the terms suppressed by this terminology.

We use the established rationalized facts
\[
 \ker(\partial|_{\mathcal P(K)})^{\operatorname{Gal}(K/\mathbb Q)}=0,
 \qquad
 \partial\mathcal P(\mathbb Q)=\Lambda^2 A(\mathbb Q)
\]
for Galois number fields $K$ \cite[\texttt{hstruct:eq:bloch-facts}]{ProveItSnapshot}. These are the Bloch-group descent and
$K_2(\mathbb Q)$ inputs already used in the manuscript. Consequently,
for an invariant $\eta$, the two kinds of reduction are equivalent,
respectively, to $\partial\eta=0$ and
$\partial\eta\in\Lambda^2 A(\mathbb Q)$.

\subsection{The existing Gram theorem and a new subfield consequence}

For $M\ge3$, set $G_M=(\mathbb Z/M\mathbb Z)^\times/\{\pm1\}$.
The all-conductor Gram theorem of the manuscript
\cite[\texttt{hstruct:thm:kernel}, \texttt{hstruct:lem:positive-Gram}]{ProveItSnapshot}
states that
\begin{equation}\label{jnew:eq:kernel}
 \sum_{a\in G_M}c_a\beta_M(a)=0
 \quad\Longleftrightarrow\quad c_a=c_{a^{-1}}\quad(a\in G_M).
\end{equation}
In particular $\beta_M(a)=0$ precisely when $a^2\equiv\pm1\pmod M$.
Its proof provides more information than this zero test. Put
\[
 v_k=(\log|1-\zeta_M^{bk}|)_{b\in G_M},\qquad
 A_M=\sum_{k=1}^{M-1}v_kv_k^{\mathsf T},\qquad
 (P_av)_b=v_{ba}.
\]
Then $A_M$ is positive definite, commutes with every $P_a$, and the
logarithmic exterior image of $\beta_M(a)$ is
\begin{equation}\label{jnew:eq:Gram-image}
 A_M(P_a-P_{a^{-1}}).
\end{equation}
The positivity uses all roots, including roots of lower order: for a
nontrivial even character of primitive conductor $f\mid M$, the
$k=M/f$ vector has a nonzero Fourier coefficient proportional to
$\tau(\chi)L(1,\overline\chi)$. Prime-order roots supply the trivial
character. This is an exact nonvanishing argument, not a numerical
determinant test. We never assume that the logarithmic exterior map
is injective on the whole exterior square.

\begin{theorem}[No proper real subfield for a nonzero individual symbol]
\label{jnew:thm:subfield}
Let $L=\mathbb Q(\zeta_M)^+$ and let $E$ be a proper subfield of $L$.
If the element $\beta_M(a)$ belongs to the image of
$\Lambda^2 A(E)$ in $\Lambda^2 A(L)$, then $\beta_M(a)=0$.
\end{theorem}
\begin{proof}
The symbols do belong to the real field, because
\[
 \langle1-\zeta_M^k\rangle
 =\tfrac12\langle(1-\zeta_M^k)(1-\zeta_M^{-k})\rangle.
\]
Let $H=\operatorname{Gal}(L/E)$, a nontrivial subgroup of $G_M$.
Logarithmic vectors of elements of $E$ lie in
$W=(\mathbb R^{G_M})^H$. A skew matrix arising from
$\Lambda^2 W$ annihilates $W^\perp$. On the complex character line
of a character $\chi$ nontrivial on $H$, formula
\eqref{jnew:eq:Gram-image} therefore gives
\[
 \lambda_\chi\bigl(\chi(a)-\chi(a)^{-1}\bigr)=0,
 \qquad\lambda_\chi>0.
\]
Thus $\chi(a)^2=1$ for every such character. Let $X_0$ be the
proper subgroup of characters trivial on $H$, and choose
$\chi_0\notin X_0$. For every $\psi\in X_0$, both $\chi_0$ and
$\chi_0\psi$ are outside $X_0$. Dividing their squared values gives
$\psi(a)^2=1$. Thus all characters of $G_M$ are trivial at $a^2$.
Characters separate points, so $a^2=1$ in $G_M$, and
\eqref{jnew:eq:kernel} gives $\beta_M(a)=0$.
\end{proof}

This statement concerns a single $\beta_M(a)$. It does not assert
that every nonzero linear combination of symbols has this property.

\subsection{Two exact dyadic identities}

\begin{lemma}[Odd conductor doubling]\label{jnew:lem:odd-double}
If $m$ is odd and $a$ is odd with $(a,m)=1$, then, inside
$\mathbb Q(\zeta_{2m})=\mathbb Q(\zeta_m)$,
\begin{equation}\label{jnew:eq:odd-double}
 \beta_{2m}(a)=3\beta_m(a)-\beta_m(2a)-\beta_m(a/2),
\end{equation}
where division by $2$ means multiplication by its inverse modulo $m$.
\end{lemma}
\begin{proof}
The even-indexed roots contribute $\beta_m(a)$. Write the odd-indexed
roots as $-\zeta_m^j$, $0\le j<m$. The $j=0$ term is
$\langle2\rangle\wedge\langle2\rangle=0$.
For $j\ne0$, put $u_j=\langle1-\zeta_m^j\rangle$. Factorization gives
$\langle1+\zeta_m^j\rangle=u_{2j}-u_j$. The remaining contribution is
\[
 \sum_{j=1}^{m-1}(u_{2aj}-u_{aj})\wedge(u_{2j}-u_j)
 =2\beta_m(a)-\beta_m(2a)-\beta_m(a/2).
\]
Adding the even-root contribution proves the identity.
\end{proof}

\begin{lemma}[Even conductor second difference]\label{jnew:lem:even-difference}
Let $e\ge2$ be even and $(a,2e)=1$. Put
\[
 \Delta_e(a)=\beta_{2e}(a)-2\beta_e(a)+\beta_{e/2}(a).
\]
Then
\begin{equation}\label{jnew:eq:Delta-zero}
 \Delta_e(a)=0
 \quad\Longleftrightarrow\quad a^2\equiv1\pmod{2e}.
\end{equation}
\end{lemma}
\begin{proof}
For $e=2$ all three symbols vanish and every odd square is $1$
modulo $4$. Suppose $e\ge4$. Let $N$ be the normalized field norm
from $\mathbb Q(\zeta_{2e})$ to $\mathbb Q(\zeta_e)$, acting on
$A$ by division by the extension degree $2$. Even-indexed roots
are fixed. If $k$ is odd, then
\[
 N\langle1-\zeta_{2e}^k\rangle
 =\tfrac12\langle1-\zeta_e^k\rangle,
\]
because the nontrivial automorphism sends $\zeta_{2e}$ to
$-\zeta_{2e}$. Every odd residue modulo $e$ has two odd lifts
modulo $2e$, hence
\begin{equation}\label{jnew:eq:dyadic-norm}
 (\Lambda^2N)\beta_{2e}(a)
 =\beta_e(a)+\tfrac12\bigl(\beta_e(a)-\beta_{e/2}(a)\bigr).
\end{equation}
It follows that
\[
 (\Lambda^2N)\Delta_e(a)
 =-\tfrac12\bigl(\beta_e(a)-\beta_{e/2}(a)\bigr).
\]
If $\Delta_e(a)=0$, then $\beta_e(a)=\beta_{e/2}(a)$, and its
definition gives $\beta_{2e}(a)=\beta_e(a)$. The right side belongs
to $\Lambda^2 A(\mathbb Q(\zeta_e)^+)$. This real field has index
$2$ in $\mathbb Q(\zeta_{2e})^+$, so Theorem
\ref{jnew:thm:subfield} implies $\beta_{2e}(a)=0$.
Thus $a^2\equiv\pm1\pmod{2e}$. Since $4\mid2e$ and $a$ is odd,
the negative sign is impossible. Conversely,
$a^2\equiv1\pmod{2e}$ makes each of the three symbols zero.
\end{proof}

\subsection{The two odd-conductor coefficient tests}

For odd $m$ and $(a,m)=1$, define
\[
 D_m(a)=\beta_m(a)-\beta_m(a/2),\qquad
 T_m(a)=\beta_m(2a)-2\beta_m(a)+\beta_m(a/2).
\]
The convention at $m=1$ is zero.

\begin{lemma}\label{jnew:lem:DT}
For every odd $m\ge1$ and every unit $a$ modulo $m$,
\begin{align}
 D_m(a)=0&\quad\Longleftrightarrow\quad m\in\{1,3,5\},
 \label{jnew:eq:D-zero}\\
 T_m(a)=0&\quad\Longleftrightarrow\quad a^2\equiv\pm1\pmod m.
 \label{jnew:eq:T-zero}
\end{align}
\end{lemma}
\begin{proof}
By \eqref{jnew:eq:kernel}, equality of two symbols
$\beta_m(a)=\beta_m(b)$ occurs only when $a=b$ in $G_m$ or when
both classes are self-inverse. Take $b=a/2$. The first alternative
gives $2=1$ in $G_m$, hence $m=1$ or $3$. In the second,
both $a^2$ and $(a/2)^2$ are signs, so $4$ is a sign modulo $m$;
thus $m=1,3,5$. Conversely, the symbol is identically zero at these
three odd conductors.

For the second assertion, use the group algebra $\mathbb R[G_m]$.
Let $g$ be the class of $2$ and let
$V_a=[a]-[a^{-1}]$. The image of $T_m(a)$ under the exact coefficient
model is
\[
 (P_g+P_{g^{-1}}-2I)V_a.
\]
On each orbit of multiplication by $g$, the kernel of this cyclic
Laplacian consists of constants. If $V_a\ne0$, its support consists
of two points with opposite coefficients. It cannot be constant
on nontrivial $g$-orbits. If $g=1$, then $m=1$ or $3$ and
$V_a=0$ anyway. Thus $T_m(a)=0$ forces $V_a=0$, which is precisely
the square congruence. The converse follows from the same displayed
formula.
\end{proof}

\subsection{Separation and classification}

If $r,s$ are coprime, normalized norms on the individual exterior
factors separate symbols from $\mathbb Q(\zeta_r)$ and
$\mathbb Q(\zeta_s)$. In detail, every summand of a $\beta$-symbol
has conjugate factors. Its normalized norm to $\mathbb Q$ has equal
factors and therefore zero wedge. A normalized norm from the compositum
to one of the two fields fixes symbols in that field and kills symbols
from the other, by linear disjointness. Consequently, for combinations
$U_r,U_s$ of $\beta$-symbols at divisors of $r,s$,
\begin{equation}\label{jnew:eq:separation}
 U_r-U_s\in\Lambda^2 A(\mathbb Q)
 \quad\Longleftrightarrow\quad U_r=U_s=0.
\end{equation}
Indeed, taking either field norm and then the norm to $\mathbb Q$
first shows that the rational wedge must be zero, and then that both
components vanish. This uses norms on factors; averaging a
Galois-invariant exterior symbol itself would not give separation.

\begin{theorem}[All rational values of the companion]
\label{jnew:thm:classification}
Let $p,q$ be coprime positive integers.
\begin{enumerate}
\item If $p$ and $q$ are odd, $J(p/q)$ has a formal
      rational-dilogarithm reduction exactly when
      $p,q\in\{1,3,5\}$. It has a formal logarithmic reduction
      exactly when $p=q=1$.
\item Otherwise, let $e$ be the even member of $\{p,q\}$ and $o$
      the odd member. Formal rational-dilogarithm and formal logarithmic
      reductions are equivalent, and occur exactly when
\begin{equation}\label{jnew:eq:main-congruences}
 \boxed{\qquad e^2\equiv\pm1\pmod o,
                  \qquad o^2\equiv1\pmod{2e}.\qquad}
\end{equation}
\end{enumerate}
Congruences modulo $1$ are vacuous.
\end{theorem}
\begin{proof}
If $p,q$ are odd, equations \eqref{jnew:eq:xi-boundary} and
\eqref{jnew:eq:odd-double} give the exact boundary
\begin{equation}\label{jnew:eq:odd-boundary}
 \partial\eta_{p/q}
 =D_q(p)-D_p(q)+\langle2\rangle\wedge\langle p/q\rangle.
\end{equation}
The field separation \eqref{jnew:eq:separation} and
\eqref{jnew:eq:D-zero} show that this boundary is rational precisely
when $p,q\in\{1,3,5\}$. The normalized norm to $\mathbb Q$ sends
the boundary to $\langle2\rangle\wedge\langle p/q\rangle$.
Since $p/q$ has odd numerator and denominator, its prime-exponent
vector is independent of that of $2$ unless $p=q=1$. Thus the
boundary is zero only in that case.

Now suppose $p=e$ is even and $q=o$ is odd. The argument $x/2$
reduces to $(e/2)/o$, and the rational wedge terms cancel exactly.
Hence
\begin{equation}\label{jnew:eq:even-boundary}
 \partial\eta_{e/o}=T_o(e)-\Delta_e(o).
\end{equation}
The fields with conductors $o$ and $2e$ are linearly disjoint.
By \eqref{jnew:eq:separation}, a rational boundary is possible
exactly when both blocks vanish; it is then zero. Lemmas
\ref{jnew:lem:even-difference} and \ref{jnew:lem:DT} give precisely
\eqref{jnew:eq:main-congruences}.

Finally, $\xi_x+\xi_{1/x}=0$ in the invariant pre-Bloch group:
their boundaries cancel, and invariant Bloch descent applies.
Pairing the three terms in \eqref{jnew:eq:eta} gives
$\eta_x+\eta_{1/x}=0$. This handles odd numerator and even
denominator. In each positive case, invariant Bloch descent also
proves the asserted converse from the boundary test to the actual
formal reduction.
\end{proof}

\begin{remark}[Why the doubled modulus cannot be dropped]
The pair $(e,o)=(8,3)$ satisfies $e^2\equiv1\pmod o$ and
$o^2\equiv1\pmod e$, but $o^2=9\not\equiv1\pmod{16}$.
Thus $F(3/8)-F(1)$ has the manuscript's rational-dilogarithm core,
while $J(3/8)$ fails the formal rational-dilogarithm test. Similarly
$(10,3)$ passes the two signed congruences for $F$ but fails the
condition for $J$. These are exact algebraic obstructions, not
failed integer-relation searches.
\end{remark}

\begin{corollary}[Two infinite families]\label{jnew:cor:families}
For every integer $n\ge1$, $J(n/(n+1))$ has a formal logarithmic
reduction. For every even integer $e\ge2$, both
$J(e/(e^2-1))$ and $J(e/(e^2+1))$ do as well.
\end{corollary}
\begin{proof}
For consecutive integers, write the odd one as $o=e\pm1$.
Then $e^2\equiv1\pmod o$ and
$o^2-1=e(e\pm2)$ is divisible by $2e$.
For $o=e^2\pm1$, the first condition is immediate and
$o^2-1=e^2(e^2\pm2)$ is divisible by $2e$ because $e$ is even.
\end{proof}
